\documentclass[11pt]{article}
\usepackage[a4paper,margin=1in]{geometry}
\usepackage[T1]{fontenc}
\usepackage{palatino}
\usepackage{enumitem}
\usepackage[hidelinks]{hyperref}
\usepackage{titlesec}
\usepackage{array}

\titlespacing*{\section}{0pt}{1.2ex plus .3ex minus .2ex}{0.6ex plus .1ex}
\titlespacing*{\subsection}{0pt}{1.0ex plus .2ex minus .2ex}{0.4ex plus .1ex}

\begin{document}

\begin{center}
{\Large\bfseries Formalising a Price-Time-Priority Matching Engine in Veil}\\[6pt]
Chen Zhikun \;\textbar\; A0310130J \;\textbar\; Solo Project\\[2pt]
{\small April 8, 2026}
\end{center}

\medskip

\section{Problem Statement}

I propose to formalise a simplified electronic limit order matching engine in
Veil. The state is a transition system: a buy-side order book sorted by
descending price, a sell-side book sorted by ascending price, a trade log,
and per-account cash and asset holdings. Orders carry an ID, account, side,
limit price, quantity, and timestamp. The transitions are \texttt{submit},
\texttt{match}, and possibly \texttt{cancel}.

The verification goal is to show that the engine preserves safety properties
expected of a price-time-priority exchange---the dominant allocation rule in
real electronic markets, where orders at each price level fill in arrival
order~\cite{harris}. Concretely, the engine should never execute a trade that
violates queue discipline, creates or destroys assets, or cannot be traced
back to two live orders. The problem sits at the intersection of relational
state modelling, ordered-data reasoning, and arithmetic invariants, which is
why it fits CS5232.

The model will not try to capture a full exchange. It covers a single asset
and limit orders only; market orders, multiple assets, and possibly
cancellation or partial fills will be omitted if they blow up the state space.
The point is a clean formal model with real correctness theorems, not market
microstructure realism.

\section{Why This Project Fits the Course}

The engine is a state machine, so Veil is the natural primary tool.
One concern worth addressing: Veil was designed for distributed protocol
verification~\cite{veilpost}, and a matching engine is strictly a centralised
sequential system. But the core Veil abstraction---transition systems over
relational state with inductive invariants---does not depend on distribution.
More concretely, the engine will be modelled as a system where multiple
external agents submit orders asynchronously; the resulting interleaving of
\texttt{submit} and \texttt{match} transitions is structurally similar to the
multi-node message-passing scenarios Veil targets. The plan is to use Veil
model checking on small instances to catch specification bugs early, then
\texttt{\#check\_invariants} for the real safety proofs.

Several of the proof obligations need Lean-level reasoning that goes beyond
what Veil's relational language handles directly: conservation of total cash
and total asset quantity, preservation of book sortedness after insertions and
deletions, and correctness of fill quantities. These will live as auxiliary
Lean lemmas.

The matching procedure also has a meaningful Velvet angle. The inner loop
(``while best bid $\geq$ best ask, execute one trade and update quantities'')
is algorithmic enough to be worth implementing in Velvet with a postcondition
tying it back to the abstract Veil spec.

What makes the invariants interesting is that they couple with each other.
Queue discipline depends on the relative ranking of all live orders on the
same side, not just a single order. Conservation links the trade log to
account balances. Partial fills, if included, add residual-quantity tracking
on top of both.

\section{Planned Formalisation}

The Veil state will track: live buy and sell orders (with a ranking relation
or sorted representation per side), each order's remaining quantity, the
executed trade log, and per-account cash and asset positions. The model covers
a single asset with a finite set of accounts. To keep things simple, the
minimum version fixes a deterministic execution-price rule instead of trying
to handle venue-specific pricing conventions.

The safety properties I plan to formalise are:

\begin{itemize}[leftmargin=1.5em]
\item \emph{stable book}: after matching completes, the remaining book is not
crossed.
\item \emph{price-time priority}: trades consume the best available orders;
within a price level, the earliest order fills first.
\item \emph{conservation}: total cash and total asset quantity across all
accounts are preserved by every step.
\item \emph{no phantom trades}: every trade log entry traces back to two orders
that were live before that step.
\item \emph{quantity consistency}: no execution uses more quantity than either
matched order has remaining.
\end{itemize}

On the Lean side, I will need helper lemmas for book sortedness (insertion,
deletion) and quantity arithmetic. If I get to the Velvet component, the plan
is a reference matching loop whose postcondition ties the output trade list
and residual book back to the Veil properties above.

The starting scope is deliberately narrow: one asset, limit orders only, no
market orders, no cancellation, and probably full fills only at first (to
avoid residual-quantity bookkeeping). Partial fills and cancellation are
stretch goals.

\section{Feasibility by April 27, 2026}

The scope can be tightened aggressively without making the result trivial:
single asset, limit orders only, safety properties only (no liveness or
throughput), and no market orders. Even with all these cuts, the engine still
has two ordered books, provenance constraints on trades, and balance
arithmetic---enough for a non-trivial proof.

I have been working with Lean throughout this course and am comfortable with
tactic-mode proofs, so the main uncertainty is not the proof infrastructure
but the state-space size for model checking. If bounded instances are too
slow, I can fall back to purely deductive proofs with
\texttt{\#check\_invariants}.

Tentative schedule:

\medskip
\noindent
\begin{tabular}{@{}>{\raggedright}p{3.2cm} p{10.4cm}@{}}
Apr 8--12  & Finalise abstract state, order representation, and basic transition relation in Veil; validate with small-instance model checking.\\[3pt]
Apr 13--18 & Formalise main safety properties and prove core invariants with \texttt{\#check\_invariants}; add Lean helper lemmas as needed.\\[3pt]
Apr 19--23 & Extend model to partial fills or cancellation if base proof is stable, or implement matching loop in Velvet.\\[3pt]
Apr 24--27 & Write report, polish development, reduce scope if necessary to ensure a clean verified core.
\end{tabular}
\medskip

The main risk is state-space explosion: two order books plus a trade log can
get expensive fast for model checking, and encoding price-time priority is
harder than plain sortedness. To control this I will start with bounded price
and quantity ranges and a small account set. If the full queue-discipline
theorem turns out to be too heavy, the fallback is to model-check it on small
instances only and prove a weaker but still useful local property
deductively---``each match step picks the current best bid and ask.''

\section{Deliverables}

\paragraph{Minimum.} A Veil specification of the single-asset matching engine;
at least one checked safety theorem combining stable-book, conservation, and
trade-provenance properties; a small finite model-checking setup; Lean helper
lemmas for sortedness and quantity arithmetic; and a short report.

\paragraph{Stretch goals (in rough priority order).}
Partial fills with correct residual-quantity updates;
cancellation with a proof that cancelled orders never appear in later trades;
a Velvet matching loop verified against an abstract postcondition;
a stronger trace-level price-time-priority theorem for the full trade log.

% Conclusion section removed to save space---the deliverables section
% already summarises the expected outcome.

\begin{thebibliography}{9}
\bibitem{cs5232proj} Ilya Sergey. ``Research Project: Guidelines and Deliverables.'' CS5232: Formal Specification and Design Techniques, 2026. \url{https://ilyasergey.net/CS5232/projects.html}
\bibitem{veilpost} Ilya Sergey. ``Verifying Distributed Protocols in Veil.'' Proofs and Intuitions, February 9, 2026. \url{https://proofsandintuitions.net/2026/02/09/distributed-verification-veil/}
\bibitem{velvetpost} Ilya Sergey. ``Multi-Modal Program Verification in Velvet.'' Proofs and Intuitions, January 21, 2026. \url{https://proofsandintuitions.net/2026/01/21/multi-modal-verification-velvet/}
\bibitem{loom} verse-lab. ``Loom: A Framework for Automated Generation of Foundational Multi-Modal Verifiers.'' GitHub repository. \url{https://github.com/verse-lab/loom}
\bibitem{harris} Larry Harris. \emph{Trading and Exchanges: Market Microstructure for Practitioners}. Oxford University Press, 2003.
\end{thebibliography}

\end{document}
